%% ============================================================
%% Chapter 7 — Error Analysis
%% ============================================================
\chapter{Error Analysis}
\label{chap:error_analysis}

\section{Error Taxonomy}
\label{sec:err_taxonomy}

We categorise errors made by the Answerability Classifier into three
orthogonal failure modes. Table~\ref{tab:error_taxonomy} provides a structured
overview; Sections~\ref{sec:err_retrieval}–\ref{sec:err_classification} discuss
each in detail.

\begin{table}[htbp]
  \centering
  \caption{Error taxonomy for the RAG-with-abstention pipeline.}
  \label{tab:error_taxonomy}
  \begin{tabular}{llp{6.5cm}}
    \toprule
    \textbf{Error Type} & \textbf{Locus} & \textbf{Description} \\
    \midrule
    Retrieval Failure      & BM25 retriever  & The ground-truth passage is not in
                                               the top-$k$ results; the reader
                                               has no access to relevant text. \\
    Extraction Failure     & RoBERTa reader  & The ground-truth passage is retrieved,
                                               but the reader extracts an incorrect
                                               or empty span. \\
    Classification Failure & Abstention layer & The classifier makes the wrong
                                               abstention decision regardless of
                                               retrieval and extraction quality. \\
    \bottomrule
  \end{tabular}
\end{table}

\subsection{Retrieval Failure}
\label{sec:err_retrieval}

Retrieval failures arise when BM25 cannot surface the relevant passage among the
top-$k$ results. This happens most often when:

\begin{itemize}
  \item The question contains synonyms or paraphrases not present in the
        passage (lexical mismatch).
  \item The question is about a rare entity that appears in many passages,
        causing relevant passages to be outranked by topically similar but
        factually different ones.
  \item The passage is very short and receives a low BM25 score due to length
        normalisation, even if it is the most relevant.
\end{itemize}

\noindent \textbf{Example.} Consider the question \textit{``What year did the
Treaty of Paris end the conflict?''} In this case, BM25 may retrieve passages
about the Treaty of Paris of 1763 (Seven Years' War), 1783 (American
Revolutionary War), or 1898 (Spanish–American War), depending on the broader
corpus context. If the ground-truth passage concerns the 1783 treaty but BM25
ranks it third, the reader operates on irrelevant text, leading to either an
empty answer or a hallucinated date.

When retrieval fails on an \emph{unanswerable} question, the system may
accidentally produce the correct abstention—but for the wrong reason.
Conversely, when retrieval fails on an \emph{answerable} question, the system
faces an extraction failure even before the abstention decision.

\subsection{Extraction Failure}
\label{sec:err_extraction}

Even when the correct passage is retrieved (top-1), the RoBERTa reader may fail
to extract the right span. Common sub-cases include:

\begin{itemize}
  \item \textbf{Off-by-one extraction.} The reader returns a span that is
        slightly wider or narrower than the ground-truth answer (e.g.,
        extracting \textit{``the 14th of July, 1789''} when the answer is
        simply \textit{``14 July 1789''}).

  \item \textbf{Plausible distractor selection.} The passage contains
        several entity mentions of the same type (e.g., multiple dates), and
        the reader selects the wrong one.

  \item \textbf{Wrong no-answer prediction.} The reader assigns a high
        \texttt{qa\_no\_answer\_probability} to a passage that actually
        contains the answer, triggering an extraction abstention even though
        the retrieval was successful.
\end{itemize}

\noindent \textbf{Example.} \textit{Question:} ``What is the boiling point of
water?'' \textit{Passage:} ``Water boils at 100°C under standard atmospheric
pressure of 101.3 kPa.'' A reader that erroneously extracts \textit{``101.3
kPa''} instead of \textit{``100°C''} commits an extraction failure.

\subsection{Classification Failure}
\label{sec:err_classification}

Classification failures occur at the abstention layer: the classifier assigns
the wrong answerability label, independent of whether retrieval and extraction
are correct. The two sub-types are:

\begin{itemize}
  \item \textbf{False Positive (FP, $n=46$).} The classifier predicts
        \emph{answerable} ($\hat{p} \geq 0.28$) for an unanswerable question.
        These are \emph{false answers}: the system returns a spurious span.
        In the test set, 46 out of 600 unanswerable questions receive a
        false answer, giving a False-Answer Rate of $46/600 \approx 7.7\%$
        among \emph{questions the classifier attempts to answer} (and 17.76\%
        of all unanswerable questions).

  \item \textbf{False Negative (FN, $n=27$).} The classifier predicts
        \emph{unanswerable} ($\hat{p} < 0.28$) for an answerable question.
        These are \emph{false abstentions}: the system withholds a correct
        answer. In the test set, 27 out of 600 answerable questions are
        suppressed, giving a False-Abstention Rate of $27/600 = 4.5\%$.
\end{itemize}

False Positives are more numerous than False Negatives (46 vs.\ 27),
consistent with the low threshold $\delta^* = 0.28$ that biases the system
toward answering. The asymmetry reflects the choice to optimise F1—which
penalises false abstentions (FN) more than false answers (FP) in this
evaluation formulation.

\section{Qualitative Case Analysis}
\label{sec:err_cases}

We examine representative examples from the three error categories.

\paragraph{Case 1 — Retrieval Failure on Answerable Question.}
\textit{Question:} [INSERT FROM NOTEBOOKS — example unanswered question due to
retrieval failure]\\
\textit{Ground-truth passage:} [INSERT FROM NOTEBOOKS]\\
\textit{Top-1 retrieved passage:} [INSERT FROM NOTEBOOKS]\\
\textit{Analysis:} The question uses domain-specific terminology not in the
ground-truth passage title, causing BM25 to rank a topically adjacent but
wrong passage first. Because the retrieved passage is irrelevant, the reader
produces an empty or incorrect span, and the classifier—seeing low retrieval
scores—correctly abstains. However, the abstention is \emph{accidental}: the
system reached the right answer for the wrong reason.

\paragraph{Case 2 — Classification Failure (False Positive).}
\textit{Question:} [INSERT FROM NOTEBOOKS — example FP case from
test\_classifier\_predictions.csv]\\
\textit{Analysis:} Despite the question being unanswerable given the corpus,
the BM25 retriever returns a passage with a relatively high score
(\texttt{top\_bm25\_score}~$\approx$~[INSERT]), and the reader produces a
confident-sounding span with \texttt{qa\_score}~$\approx$~[INSERT]. These
signals mislead the classifier into predicting \emph{answerable}, producing a
false answer.

\paragraph{Case 3 — Classification Failure (False Negative).}
\textit{Question:} [INSERT FROM NOTEBOOKS — example FN case]\\
\textit{Analysis:} The correct passage is retrieved (rank 1), and the reader
extracts the right span. However, the \texttt{qa\_no\_answer\_probability} is
anomalously elevated (e.g., because the answer span is a numeric entity with
unusual token characteristics), and the classifier—weighting this feature
heavily—predicts \emph{unanswerable}, causing an unnecessary abstention.

\section{Summary}
\label{sec:err_summary}

The error analysis reveals that:

\begin{enumerate}
  \item Retrieval failures constitute a fundamental bottleneck. When BM25
        cannot retrieve the right passage, neither the reader nor the
        abstention layer can compensate.

  \item Classification failures (FP and FN) account for 73 out of 1,200 test
        questions (6.1\% of total), suggesting substantial room for improvement
        with richer features or a non-linear classifier.

  \item The classifier's main advantage over the BM25 Threshold strategy is
        its ability to integrate reader signals (\texttt{qa\_score},
        \texttt{qa\_no\_answer\_probability}) that capture post-retrieval
        evidence unavailable to a pure retrieval-based threshold.
\end{enumerate}
